\documentclass{article}
\usepackage{fuzz}

\begin{document}

\title{vox-dvri: Record-Store Path Containment}
\author{rmh}
\maketitle

% The daemon owns a recordings store rooted at a single directory it created
% 0700. A wire client names a file only by a candidate: a "kind" flag
% (absolute or relative) plus a sequence of path components. The security
% invariant (vox-zu39, P1) is that the daemon never writes, plays, or reads a
% path outside that root. This spec models candidate resolution and the store
% state so the containment invariant is stated in the schema predicate and each
% operation is shown to preserve it.

% Path components (a bare filename is a one-element sequence) and resolved
% filesystem paths.
\begin{zed}
  [COMP, PATH]
\end{zed}

% A candidate is absolute or relative. Relative is the only acceptable kind:
% voxd's cwd is not the caller's shell, and an absolute candidate is a direct
% escape from the root.
\begin{zed}
  Flag ::= isAbs | isRel
\end{zed}

% The distinguished traversal component ".."; its presence lets a relative
% candidate climb out of the root, so it is unsafe.
\begin{axdef}
  dotdot: COMP
\end{axdef}

% The store root, an abstract "contained under" relation on paths (reflexive on
% the root itself), and the resolver that maps a candidate to the path it would
% land at.
\begin{axdef}
  root: PATH \\
  under: PATH \rel PATH \\
  resolve: (Flag \cross \seq COMP) \pfun PATH
\end{axdef}

% A candidate is SAFE exactly when it is relative, non-empty, and free of any
% ".." component. This is the wire-input rule the daemon enforces before any
% filesystem touch (it also rejects separators/NUL, modelled here as structural
% illegality that keeps a name to safe components). The second predicate is the
% resolver's containment guarantee: every SAFE candidate resolves to a concrete
% path contained under the root -- the property the implementation's
% ``(root / name).resolve()`` plus ``is_relative_to(root)`` check must satisfy.
\begin{axdef}
  safe: \power (Flag \cross \seq COMP)
\where
  \forall k: Flag; s: \seq COMP @ \\
  \t1 (k, s) \in safe \iff \\
  \t2 (k = isRel \land dotdot \notin \ran s \land s \neq \langle \rangle) \\
  \forall k: Flag; s: \seq COMP @ \\
  \t1 (k, s) \in safe \implies \\
  \t2 (k, s) \in \dom resolve \land (resolve~(k, s), root) \in under
\end{axdef}

% The store: the set of paths that hold stored recordings. The invariant --
% ``every stored write stays within the root'' -- is carried in the predicate,
% not in prose.
\begin{schema}{Store}
  stored: \power PATH
\where
  \forall p: stored @ (p, root) \in under
\end{schema}

\begin{schema}{InitStore}
  Store'
\where
  stored' = \emptyset
\end{schema}

% Place a recording named by a candidate. Precondition: the candidate is SAFE
% (an absolute, traversing, or empty candidate never reaches here -- see
% Reject). The landed path is the resolved path, which the containment
% guarantee proves is under the root, so the Store invariant is preserved.
\begin{schema}{Place}
  \Delta Store \\
  kind?: Flag \\
  comps?: \seq COMP \\
  path!: PATH
\where
  (kind?, comps?) \in safe \\
  path! = resolve~(kind?, comps?) \\
  stored' = stored \cup \{ path! \}
\end{schema}

% Report codes for the rejection path.
\begin{zed}
  REPORT ::= accepted | rejected
\end{zed}

% Reject an unsafe candidate: absolute, or containing "..", or empty. The store
% is unchanged -- no write escapes the root because no write happens.
\begin{schema}{Reject}
  \Xi Store \\
  kind?: Flag \\
  comps?: \seq COMP \\
  report!: REPORT
\where
  (kind?, comps?) \notin safe \\
  report! = rejected
\end{schema}

% Play and Fetch both address an existing store recording by its resolved path.
% Precondition: the reference is already in ``stored'' -- and every member of
% ``stored'' is under the root by the invariant -- so neither operation can
% touch a path outside the sandbox. Read-only: the store is unchanged.
\begin{schema}{Access}
  \Xi Store \\
  ref?: PATH \\
  report!: REPORT
\where
  ref? \in stored \\
  report! = accepted
\end{schema}

\begin{zed}
  Play == Access \\
  Fetch == Access
\end{zed}

% A reference that is not a stored recording (an absolute path, a traversal, or
% simply an unknown name that resolved outside ``stored'') is refused.
\begin{schema}{AccessReject}
  \Xi Store \\
  ref?: PATH \\
  report!: REPORT
\where
  ref? \notin stored \\
  report! = rejected
\end{schema}

\end{document}
